\subsection{Comparison with published clustering rules}\label{sec:rules}

\tabref{tab:main} reports every method on 18 instances across
3 problem sizes, at $\varphi = 100$, two alignment
levels and three seeds each. Its cost columns are means over the 18
instances, and its \emph{vs.\ ours} columns are paired on the per-instance
differences instead, since the instances differ in difficulty by more than the
methods differ from one another.

\begin{table}[t]\centering
\sbox{\tabbox}{\small\begin{tabular}{@{}l rrrrr rrrrr@{}}
\toprule
& \multicolumn{5}{c}{as published} & \multicolumn{5}{c}{+ descent}\\
\cmidrule(lr){2-6}\cmidrule(lr){7-11}
method & travel & tard. & cour. & cost & vs.\ ours & travel & tard. & cour. & cost & vs.\ ours\\
\midrule
\textsc{sweep} & 2163 & 833 & 22.7 & 5263 & $-29.3\%$ & 1932 & 35 & 22.7 & 4234 & $-13.3\%$\\
\textsc{fisher} & 2487 & 578 & 34.0 & 6464 & $-42.6\%$ & 1692 & 12 & 26.6 & 4359 & $-16.1\%$\\
$k$-means & 2174 & 410 & 34.0 & 5984 & $-38.8\%$ & 1876 & 18 & 28.1 & 4700 & $-21.6\%$\\
spatio-temporal $k$-means & 2217 & 800 & 34.0 & 6417 & $-43.0\%$ & 1932 & 31 & 28.6 & 4819 & $-22.7\%$\\
capacitated $k$-means & 2284 & 449 & 34.0 & 6133 & $-40.0\%$ & 1922 & 12 & 28.4 & 4773 & $-22.5\%$\\
Ward agglomerative & 2176 & 326 & 34.0 & 5902 & $-36.8\%$ & 1822 & 40 & 28.2 & 4684 & $-19.5\%$\\
\addlinespace[2pt]
best of six & 2149 & 809 & 22.9 & 5247 & $-28.8\%$ & 1779 & 18 & 23.9 & 4186 & $-12.1\%$\\
\midrule
\textbf{this paper} & \multicolumn{5}{c}{} & 1206 & 105 & 23.2 & \textbf{3628} & ---\\
ALNS & \multicolumn{5}{c}{} & 2040 & 92 & 23.1 & 4438 & $-14.2\%$\\
\bottomrule\end{tabular}}%
\setlength{\tabw}{1.00\textwidth}%
\ifdim\wd\tabbox>\tabw\setlength{\tabw}{\wd\tabbox}\fi
\begin{minipage}{\tabw}
\centering
\caption{$|\mathcal{R}| = 40, 100, 200$ requests, $|\mathcal{K}| = 12, 30, 60$ couriers, $Q_k = 5$, $\varphi = 100$, $\rho = 1$. Cost and its components, means over all instances. \emph{As published} uses the rule in its original form; \emph{+ descent} adds the improvement step of Section~\ref{sec:descent}. \emph{vs.\ ours} is the mean paired percentage cost difference against the proposed method, negative favouring it; each is a Wilcoxon signed-rank test over the same instances, and every one of them has $p<10^{-3}$. \emph{Best of six} takes the cheapest rule on each instance separately, so it is a stricter baseline than any single row. The proposed method and ALNS have no \emph{as published} form, so they appear once. Cost is travel${}+{}$tardiness${}+\varphi\times{}$couriers.}
\label{tab:main}
\small\begin{tabular*}{\tabw}{@{\extracolsep{\fill}}l rrrrr rrrrr@{}}
\toprule
& \multicolumn{5}{c}{as published} & \multicolumn{5}{c}{+ descent}\\
\cmidrule(lr){2-6}\cmidrule(lr){7-11}
method & travel & tard. & cour. & cost & vs.\ ours & travel & tard. & cour. & cost & vs.\ ours\\
\midrule
\textsc{sweep} & 2163 & 833 & 22.7 & 5263 & $-29.3\%$ & 1932 & 35 & 22.7 & 4234 & $-13.3\%$\\
\textsc{fisher} & 2487 & 578 & 34.0 & 6464 & $-42.6\%$ & 1692 & 12 & 26.6 & 4359 & $-16.1\%$\\
$k$-means & 2174 & 410 & 34.0 & 5984 & $-38.8\%$ & 1876 & 18 & 28.1 & 4700 & $-21.6\%$\\
spatio-temporal $k$-means & 2217 & 800 & 34.0 & 6417 & $-43.0\%$ & 1932 & 31 & 28.6 & 4819 & $-22.7\%$\\
capacitated $k$-means & 2284 & 449 & 34.0 & 6133 & $-40.0\%$ & 1922 & 12 & 28.4 & 4773 & $-22.5\%$\\
Ward agglomerative & 2176 & 326 & 34.0 & 5902 & $-36.8\%$ & 1822 & 40 & 28.2 & 4684 & $-19.5\%$\\
\addlinespace[2pt]
best of six & 2149 & 809 & 22.9 & 5247 & $-28.8\%$ & 1779 & 18 & 23.9 & 4186 & $-12.1\%$\\
\midrule
\textbf{this paper} & \multicolumn{5}{c}{} & 1206 & 105 & 23.2 & \textbf{3628} & ---\\
ALNS & \multicolumn{5}{c}{} & 2040 & 92 & 23.1 & 4438 & $-14.2\%$\\
\bottomrule\end{tabular*}
\end{minipage}
\end{table}

The \emph{vs.\ ours} columns of \tabref{tab:main} pair the proposed method
against each of the six rules separately and against the best of them. The
comparison we lead with is the last, and it is the strictest: on each instance
the baseline is whichever of the six is cheapest \emph{on that instance}, so
the margin against it is necessarily smaller than the margin against any rule
taken alone, as the spread of those rows shows. Against that baseline our proposed method is
29\% cheaper with the six as published, winning 18
of 18 instances. Giving the six our own descent narrows that to
12.1\% ($p<10^{-3}$, still 18/18). The two
figures measure different things. The
16.7 points between them are what a generic
improvement step recovers for an existing rule. The
12.1\% that survives it is the contribution of
grouping on the marginal cost of~\eqref{eq:marginal}, measured with the
improvement step common to both sides, and it is the figure the rest of this
section builds on.


The cost difference behind that 12.1\% is the sum of three
terms of the objective~\eqref{eq:globalobj} that move in different directions.
Pooled over the 18 instances and measured in the units of that objective,
our proposed method saves $573$ in travel, gives up $87$ in
tardiness and opens $0.72$ fewer couriers on average,
worth $72$ at $\varphi=100$; the three account exactly
for the $558$ it saves in cost. The advantage is therefore a
large reduction in travel bought with a small concession in lateness, and not a
uniform improvement across the objective. The margin also grows with instance size, and
\tabref{tab:bysize} gives it by size against both baselines. Against the six
strengthened rules it is $-8.9\%$ at $|\mathcal{R}|=40$, $-13.1\%$ at $|\mathcal{R}|=100$, $-14.4\%$ at $|\mathcal{R}|=200$.

\begin{table}[t]\centering
\sbox{\tabbox}{\small\begin{tabular}{@{}r rrr rrr@{}}\toprule
& \multicolumn{3}{c}{vs.\ best of six $+$ descent} & \multicolumn{3}{c}{vs.\ ALNS}\\
\cmidrule(lr){2-4}\cmidrule(lr){5-7}
$|\mathcal{R}|$ & travel & cour. & cost & travel & cour. & cost\\\midrule
$40$ & $-9.3\%$ & $-8.9\%$ & $-8.9\%$ & $-3.6\%$ & $-3.7\%$ & $-4.1\%$\\
$100$ & $-32.9\%$ & $+1.1\%$ & $-13.1\%$ & $-37.4\%$ & $+0.9\%$ & $-17.6\%$\\
$200$ & $-36.4\%$ & $-3.0\%$ & $-14.4\%$ & $-47.1\%$ & $+1.3\%$ & $-20.8\%$\\
\bottomrule\end{tabular}}%
\setlength{\tabw}{0.80\textwidth}%
\ifdim\wd\tabbox>\tabw\setlength{\tabw}{\wd\tabbox}\fi
\begin{minipage}{\tabw}
\centering
\caption{$|\mathcal{R}| = 40, 100, 200$ requests, $|\mathcal{K}| = 12, 30, 60$ couriers, $Q_k = 5$, $\varphi = 100$, $\rho = 1$. Paired differences by problem size, on the instances of \tabref{tab:main}. Each entry is the mean of the per-instance percentage differences at that size, negative favouring the proposed method. \emph{Best of six} is the cheapest of the six rules carrying our descent on that instance. Six instances at each size, two alignment levels by three seeds. A tardiness column is omitted because both baselines are close to punctual here, so a ratio against their tardiness is set by the instances where it is almost zero rather than by the difference; the pooled difference in the units of the objective is given in the text.}
\label{tab:bysize}
\small\begin{tabular*}{\tabw}{@{\extracolsep{\fill}}r rrr rrr@{}}\toprule
& \multicolumn{3}{c}{vs.\ best of six $+$ descent} & \multicolumn{3}{c}{vs.\ ALNS}\\
\cmidrule(lr){2-4}\cmidrule(lr){5-7}
$|\mathcal{R}|$ & travel & cour. & cost & travel & cour. & cost\\\midrule
$40$ & $-9.3\%$ & $-8.9\%$ & $-8.9\%$ & $-3.6\%$ & $-3.7\%$ & $-4.1\%$\\
$100$ & $-32.9\%$ & $+1.1\%$ & $-13.1\%$ & $-37.4\%$ & $+0.9\%$ & $-17.6\%$\\
$200$ & $-36.4\%$ & $-3.0\%$ & $-14.4\%$ & $-47.1\%$ & $+1.3\%$ & $-20.8\%$\\
\bottomrule\end{tabular*}
\end{minipage}
\end{table}
\subsection{The grouping cost examined}\label{sec:groupingcost}\titledParagraph{Non-additivity of the routing cost.}\label{sec:interaction}
Section~\ref{sec:guarantees} identifies the pairwise interaction of
$J(\cdot,k)$ as the exact measure of the additivity that the assignment program
assumes and routing does not have. This subsection measures it. We sampled 12,000 second differences
$I(\mathcal{S};r,r')=J(\mathcal{S}\cup\{r,r'\})-J(\mathcal{S}\cup\{r\})
-J(\mathcal{S}\cup\{r'\})+J(\mathcal{S})$ with the exact router, in which
$\mathcal{S}$ is a set of requests a courier already holds and $r,r'$ are two
further requests it might be given.
\tabref{tab:interaction} reports them by
$|\mathcal{S}|$, which runs from $0$ to $Q_k-2=3$: all of
$\mathcal{S}\cup\{r,r'\}$ has to fit one courier, so a base set of
4 leaves no room for a pair.

\begin{table}[t]\centering
\sbox{\tabbox}{\small\begin{tabular}{@{}r@{\;\;} r@{\;\;} r@{\;\;} r@{\;\;} r@{}}\toprule
$|\mathcal{S}|$ & samples & zero & median & 95th pct.\\\midrule
$0$ & 3,054 & $0\%$ & $0.538$ & $0.655$\\
$1$ & 3,068 & $27\%$ & $0.201$ & $0.462$\\
$2$ & 2,944 & $39\%$ & $0.074$ & $0.344$\\
$3$ & 2,934 & $45\%$ & $0.024$ & $0.261$\\
\bottomrule\end{tabular}}%
\setlength{\tabw}{0.80\textwidth}%
\ifdim\wd\tabbox>\tabw\setlength{\tabw}{\wd\tabbox}\fi
\begin{minipage}{\tabw}
\centering
\caption{$|\mathcal{R}| = 40$ requests, $|\mathcal{K}| = 12$ couriers, $Q_k = 5$, $\varphi = 100$, $\rho = 1$. The interaction $I(\mathcal{S};r,r')$ of~\eqref{eq:interaction}, computed with the exact router on sampled triples $(\mathcal{S},r,r')$ and grouped by $|\mathcal{S}|$, the number of requests the courier already holds. $I$ is zero for every $\mathcal{S}$, $r$ and $r'$ exactly when a courier's routing cost is additive across its requests, which is what the assignment program assumes. \emph{Zero} is the fraction of samples whose $I$ vanishes exactly; \emph{median} and \emph{95th pct.} give $|I|$ relative to the marginal it perturbs.}
\label{tab:interaction}
\small\begin{tabular*}{\tabw}{@{\extracolsep{\fill}}r@{\;\;} r@{\;\;} r@{\;\;} r@{\;\;} r@{}}\toprule
$|\mathcal{S}|$ & samples & zero & median & 95th pct.\\\midrule
$0$ & 3,054 & $0\%$ & $0.538$ & $0.655$\\
$1$ & 3,068 & $27\%$ & $0.201$ & $0.462$\\
$2$ & 2,944 & $39\%$ & $0.074$ & $0.344$\\
$3$ & 2,934 & $45\%$ & $0.024$ & $0.261$\\
\bottomrule\end{tabular*}
\end{minipage}
\end{table}

The routing cost is not additive across requests, and the table says by how
much. The typical departure from additivity is small, median
7.3 time units against a mean of 171 for
serving one of these sets, and its sign is close to balanced,
35\% negative against 38\% positive, so the errors a courier
accumulates over its requests partly cancel rather than compound. The mean
relative interaction is $-0.9\%$. In
27\% of samples it vanishes exactly, which happens when the two requests
never compete for a position in the route: each is inserted where the other
leaves the schedule untouched, so what one adds does not depend on whether the
other is there. The worst case is different: the largest interaction sampled is
307 time units, so Section~\ref{sec:guarantees} claims nothing about how
accurately any one courier's requests are costed.

Where the interaction is measured matters more than its average size, and that
settles a design choice we had made on intuition. It falls sharply with the size
of the set the marginal is measured in: at $|\mathcal{S}|=0$ no sample vanishes
and the median interaction is $0.54$ of the marginal it
perturbs, while at $|\mathcal{S}|=3$ nearly half
vanish and the median is $0.024$, a factor of
23 apart. A
marginal measured against a courier's \emph{current} workload is therefore far
more accurate than one measured against an empty courier. That is
what~\eqref{eq:marginal} does, and it is also why the outer loop recomputes
$C_{rk}$ at every pass: the set of requests a courier holds changes from one
pass to the next, and a marginal measured against requests the courier no longer
holds is measured in the wrong place.

\titledParagraph{Temporal criteria in the grouping cost.}\label{sec:temporal}\label{sec:slack}
In Section~\ref{sec:cost} we argued that a temporal coordinate is the wrong
thing to group by, because it draws together the requests that compete for a
courier rather than those it can serve in turn. That claim can be tested without
reference to our proposed method at all, because two of the six rules differ
in exactly
one respect: $k$-means groups requests by their position, and spatio-temporal
$k$-means adds a scaled time coordinate to the same algorithm. If grouping requests that fall due together
helped, the second would be the better of the two. It is not, and where it
fails is exactly where the argument predicts, as
\tabref{tab:temporal} shows.

\begin{table}[t]\centering
\sbox{\tabbox}{\small\begin{tabular}{@{}l rr r rr@{}}\toprule
& \multicolumn{2}{c}{tardiness} & & \multicolumn{2}{c}{cost}\\
\cmidrule(lr){2-3}\cmidrule(lr){5-6}
instances & $k$-m. & s.-t. & $p$ & $k$-m. & s.-t.\\\midrule
aligned & 575 & 743 & $0.164$ & 6038 & 6251\\
misaligned & 245 & 857 & $0.004$ & 5930 & 6583\\
\bottomrule\end{tabular}}%
\setlength{\tabw}{0.80\textwidth}%
\ifdim\wd\tabbox>\tabw\setlength{\tabw}{\wd\tabbox}\fi
\begin{minipage}{\tabw}
\centering
\caption{$|\mathcal{R}| = 40, 100, 200$ requests, $|\mathcal{K}| = 12, 30, 60$ couriers, $Q_k = 5$, $\varphi = 100$, $\rho = 1$. $k$-means grouping with and without a time coordinate, paired on identical instances. \emph{$k$-m.}\ is $k$-means on request midpoints; \emph{s.-t.}\ is spatio-temporal $k$-means, the same algorithm carrying the extra time coordinate of Section~\ref{sec:setup} at temporal weight $\omega=1$, so the two differ in that coordinate alone. Misaligned instances are those where spatial and temporal structure disagree, so the time coordinate carries information the spatial one does not.}
\label{tab:temporal}
\small\begin{tabular*}{\tabw}{@{\extracolsep{\fill}}l rr r rr@{}}\toprule
& \multicolumn{2}{c}{tardiness} & & \multicolumn{2}{c}{cost}\\
\cmidrule(lr){2-3}\cmidrule(lr){5-6}
instances & $k$-m. & s.-t. & $p$ & $k$-m. & s.-t.\\\midrule
aligned & 575 & 743 & $0.164$ & 6038 & 6251\\
misaligned & 245 & 857 & $0.004$ & 5930 & 6583\\
\bottomrule\end{tabular*}
\end{minipage}
\end{table}

On instances where space and time disagree, adding the time coordinate
multiplies tardiness by $3.5$ ($p=0.004$) and raises total cost by
$11.4\%$ ($p=0.008$). On instances where the two agree
the effect is not significant ($p=0.164$), which is what one would
expect: there the time coordinate merely repeats what the spatial one already
says, so it can do little harm and no good. The published remedy for time
windows in a clustering rule therefore makes the rule worse precisely when it is
supposed to help, and the reason is the one given in
Section~\ref{sec:cost}: a courier serves its group in sequence, so requests
that fall due together contend for it.


The time coordinate carries a weight, so the natural objection is that the
weight was set badly. \figrefp{fig:sensitivity}{c} sweeps the
temporal weight $\omega$ of Section~\ref{sec:setup} over 8 values. No
setting helps by more than noise. The best of them in the rule's published form
is the smallest weight tried, $\omega=0.125$, at
$-1.0\%$ against $\omega=0$ on
7 of 12 instances
($p=0.074$). From $\omega=1$ upward every weight is
significantly worse, and the cost climbs to $+18\%$ at
$\omega=8$ with tardiness climbing alongside it, from
286 to 871.
Carried by our descent the same rule is flat to within 4\%, which is
the descent repairing the grouping rather than the weight becoming useful.


That was the crude way of putting time into a grouping cost. The sophisticated one
concedes what Section~\ref{sec:cost} argues, that minimising a temporal distance
groups the requests that compete for a courier, and introduces time another
way: it
charges each pair of requests the courier cannot chain, by
$g_{\mathrm{sum}}$ or $g_{\min}$ of~\eqref{eq:chaincharge}. That is the
criterion most likely to be proposed in place of the marginal cost, so it is the
one worth measuring. Like the marginal cost of~\eqref{eq:marginal}, this is a
grouping cost, and it is used in its place, so the comparison is between two
answers to the same question. It is built and measured exactly as the six
published rules are, on the same 12 instances.

The chaining slack of Section~\ref{sec:cost} is directional.
$\operatorname{slack}(r\!\to\!r')$ is the spare time left if the courier
serves $r$ and then $r'$, and $\operatorname{slack}(r'\!\to\!r)$ is a
value computed in the other direction. A grouping cost charges a pair rather than an ordered pair, so
the two directions have to be combined into one, and there are two ways to do it.
Summing them charges a pair when \emph{either} order fails; taking their
minimum charges it only when \emph{both} do. The second is the correct reading,
since the courier chooses the best order, and both are measured below.

The two symmetrisations are not variants of one quantity, and correlating each
against the temporal separation $|\tau_r-\tau_{r'}|$ of
Section~\ref{sec:cost} shows why. Taken over all pairs of requests in these
instances, $g_{\mathrm{sum}}$ of~\eqref{eq:chaincharge} correlates $+0.95$ with
separation: it is, to a good approximation, a measure of how far apart in time
two requests are, so minimising it groups requests that fall due together, which
is what Section~\ref{sec:cost} argues against. $g_{\min}$ correlates
$-0.74$, of the opposite sign, scoring requests far apart in
time as the compatible ones. The correction does what it claims.
\tabref{tab:slack} reports what each symmetrisation
costs.

\begin{table}[t]\centering
\sbox{\tabbox}{\small\begin{tabular}{@{}l@{\;} r@{\;\;}r@{\;\;}r@{\;\;} r@{\;\;} r@{}}\toprule
grouping cost & travel & tard. & cour. & cost & vs.\ ours\\\midrule
slack, sum & 2221 & 1387 & 15.2 & 5133 & $+114.8\%$\\
slack, min & 2218 & 1673 & 16.7 & 5558 & $+134.3\%$\\
slack, sum $+$ desc. & 1291 & 11 & 14.7 & 2768 & $+17.2\%$\\
slack, min $+$ desc. & 1245 & 6 & 15.2 & 2776 & $+16.5\%$\\
best published $+$ desc. & 1184 & 13 & 14.7 & 2664 & $+12.6\%$\\
\midrule
marginal cost (ours) & 873 & 42 & 14.3 & 2349 & ---\\
\bottomrule\end{tabular}}%
\setlength{\tabw}{0.80\textwidth}%
\ifdim\wd\tabbox>\tabw\setlength{\tabw}{\wd\tabbox}\fi
\begin{minipage}{\tabw}
\centering
\caption{$|\mathcal{R}| = 40, 100$ requests, $|\mathcal{K}| = 12, 30$ couriers, $Q_k = 5$, $\varphi = 100$, $\rho = 1$. Grouping by pairwise chaining slack, in both symmetrisations, against the marginal cost. \emph{Sum} charges both orders of a pair, \emph{min} only the better one. Each is reported as published and with our descent, and all are routed by the same exact router.}
\label{tab:slack}
\small\begin{tabular*}{\tabw}{@{\extracolsep{\fill}}l@{\;} r@{\;\;}r@{\;\;}r@{\;\;} r@{\;\;} r@{}}\toprule
grouping cost & travel & tard. & cour. & cost & vs.\ ours\\\midrule
slack, sum & 2221 & 1387 & 15.2 & 5133 & $+114.8\%$\\
slack, min & 2218 & 1673 & 16.7 & 5558 & $+134.3\%$\\
slack, sum $+$ desc. & 1291 & 11 & 14.7 & 2768 & $+17.2\%$\\
slack, min $+$ desc. & 1245 & 6 & 15.2 & 2776 & $+16.5\%$\\
best published $+$ desc. & 1184 & 13 & 14.7 & 2664 & $+12.6\%$\\
\midrule
marginal cost (ours) & 873 & 42 & 14.3 & 2349 & ---\\
\bottomrule\end{tabular*}
\end{minipage}
\end{table}

Nonetheless, correcting the symmetrisation does not rescue the criterion. As published the
minimum is if anything the worse of the two
($+134.3\%$ against ours, versus
$+114.8\%$ for the sum), and once
both are given our descent they are statistically indistinguishable from each
other ($-0.5\%$, $p=0.733$) and both remain worse than
the marginal cost by $+17.2\%$ and $+16.5\%$
($p<10^{-3}$). They are also worse than the best published rule carrying
the same descent, so neither reaches the rules the idea was meant to improve on.

The components say where the difference sits, and it is not punctuality. Carrying
our descent, the slack rules are the \emph{more} punctual of the two, with
tardiness 6 and 11 against our 42, but they travel
1245 and 1291 against our 873 and open more couriers. At
$\rho=1$ that trade runs about ten to one against them: the better of the two
saves 36 units of lateness and pays 372 units of
extra distance for it. Section~\ref{sec:rho} shows the same trade against the
published rules, across a twentyfold range in $\rho$.

Getting the sign right is therefore not enough, and that is a separate point from
the correlation. A charge on pairs
asks whether two requests can follow one another; what a courier is actually
charged is the cost of a \emph{path} through all of them, and no sum over pairs
recovers it, because what a pair costs depends on what else the courier is
carrying and where the path has reached. The failure is not in how the two
directions were combined, which the minimum repairs, but in charging pairs at
all. That is the quantity \eqref{eq:marginal} evaluates directly.

In the proposed method time does not enter the grouping cost directly, as it does
in both criteria above. It enters through the marginal cost
of~\eqref{eq:marginal}, which prices a request by the lateness the courier would
actually incur on the route it would then take, at the rate $\rho$ the objective
sets. Our proposed method therefore spends on punctuality exactly what the
objective says it is worth. That is why it is the
less punctual of the two here and still the cheaper, and why the same mechanism
reverses when lateness is priced higher: Section~\ref{sec:rho} follows it across
a twentyfold range in $\rho$. A criterion with no $\rho$ in it cannot make that
trade in either direction.

\subsection{Sensitivity to instance conditions}\label{sec:sensitivity}
\titledParagraph{Spatio-temporal misalignment.}\label{sec:alignaxis}
The comparison so far uses the two ends of the alignment axis, $\alpha=1$ and
$\alpha=0$. Misalignment is measured on each instance rather than set by
$\alpha$, so intermediate values place instances along the axis and the margin
can be read as a function of it. \figrefp{fig:sensitivity}{b} sweeps
$\alpha$ over 5 levels with 8
seeds each.


The margin over the strengthened rules grows monotonically with realised
misalignment: the Spearman rank correlation between the two is
$-0.90$ ($p=0.037$), negative because a wider margin
is a more negative cost difference. Our proposed method is cheaper on every
instance at every level ($p=0.008$ at each). The bars on the right axis of the same panel
show the effect without reference to our proposed method, since they compare
only the
two baselines: the tardiness spatio-temporal $k$-means pays for its time
coordinate, relative to plain $k$-means, falls from $3.8\times$ at the
misaligned end to $1.2\times$ where space and time agree. That
ordering is what Section~\ref{sec:cost} predicts. Where the time axis repeats
what the spatial one already says, adding it costs little; where it carries
information the spatial axis does not, minimising distance along it concentrates
the requests that compete for a courier.

One value from this sweep is used later. The adapted Li \& Lim batches sit at a
realised misalignment of $0.61$, beside our $\alpha=0$ setting at the hard end
of the axis, so the smaller margin Section~\ref{sec:lilim} reports on them is
not explained by their being temporally easy.

\titledParagraph{The penalty rate.}\label{sec:rho}
In Section~\ref{sec:cost} we argued that a courier visits its requests one after
another, so two requests whose windows coincide compete for the same
courier-minute, and that grouping by temporal proximity therefore collects
exactly those competing pairs onto one courier. Grouping on the marginal cost
of~\eqref{eq:marginal} instead of on temporal proximity might therefore be
expected to buy punctuality. However, what it buys is distance: travel falls and
lateness does not. One explanation would be
that $\rho=1$ makes lateness too cheap to be worth buying, in which case pricing
it higher should close the gap, so we swept $\rho$ over a twentyfold range. It enters every part of the
method: the exact router sequences with it, the marginal costs
of~\eqref{eq:marginal} are computed with it, and the descent decides with it.
\tabref{tab:rho} reports the sweep, and
\figrefp{fig:sensitivity}{d} plots it beside the tardiness of both methods.

\begin{table}[t]\centering
\sbox{\tabbox}{\small\begin{tabular}{@{}r@{\;\;} r@{\;\;}r@{\;\;} r@{\;\;}r@{\;\;} r@{\;\;} r@{}}\toprule
& \multicolumn{2}{c}{tardiness} & \multicolumn{2}{c}{travel} & & \\
\cmidrule(lr){2-3}\cmidrule(lr){4-5}
$\rho$ & ours & base & ours & base & cost & $p$\\\midrule
$0.5$ & 65.66 & 29.41 & 564 & 696 & $-11.1\%$ & $0.008$\\
$1$ & 14.34 & 7.45 & 606 & 694 & $-9.1\%$ & $0.008$\\
$2$ & 4.46 & 2.82 & 645 & 739 & $-8.7\%$ & $0.008$\\
$5$ & 0.90 & 0.18 & 637 & 740 & $-9.3\%$ & $0.008$\\
$10$ & 0.73 & 0.03 & 644 & 720 & $-8.2\%$ & $0.008$\\
\bottomrule\end{tabular}}%
\setlength{\tabw}{0.80\textwidth}%
\ifdim\wd\tabbox>\tabw\setlength{\tabw}{\wd\tabbox}\fi
\begin{minipage}{\tabw}
\centering
\caption{$|\mathcal{R}| = 40$ requests, $|\mathcal{K}| = 12$ couriers, $Q_k = 5$, $\varphi = 100$, $\rho = 1$. Sweeping the penalty rate $\rho$, which prices earliness and tardiness alike in~\eqref{eq:globalobj}. \emph{Base} is the best of the six rules carrying our descent on each instance. Tardiness is a mean over the instances, so a small value is an average of instances that are exactly punctual and one or two that are not, rather than a small lateness everywhere. The proposed method responds to $\rho$ as it should, and its advantage stays travel at every $\rho$.}
\label{tab:rho}
\small\begin{tabular*}{\tabw}{@{\extracolsep{\fill}}r@{\;\;} r@{\;\;}r@{\;\;} r@{\;\;}r@{\;\;} r@{\;\;} r@{}}\toprule
& \multicolumn{2}{c}{tardiness} & \multicolumn{2}{c}{travel} & & \\
\cmidrule(lr){2-3}\cmidrule(lr){4-5}
$\rho$ & ours & base & ours & base & cost & $p$\\\midrule
$0.5$ & 65.66 & 29.41 & 564 & 696 & $-11.1\%$ & $0.008$\\
$1$ & 14.34 & 7.45 & 606 & 694 & $-9.1\%$ & $0.008$\\
$2$ & 4.46 & 2.82 & 645 & 739 & $-8.7\%$ & $0.008$\\
$5$ & 0.90 & 0.18 & 637 & 740 & $-9.3\%$ & $0.008$\\
$10$ & 0.73 & 0.03 & 644 & 720 & $-8.2\%$ & $0.008$\\
\bottomrule\end{tabular*}
\end{minipage}
\end{table}

As $\rho$ rises from 0.5 to 10, the mean tardiness of our proposed method
falls from 65.7 to 0.7 and its mean travel rises from 564 to
644.
The comparison with the baselines does not move with it. At every $\rho$ our
solutions are the later of the two, by roughly a factor of two while lateness is
still material and by less than a unit once $\rho$ has driven both near zero, and
at every $\rho$ our travel is lower, by $76$ to $132$ units.
The cost advantage is flat across the range, $-11.1\%$ at
$\rho=0.5$ and $-8.2\%$ at $\rho=10$, on every instance at
every $\rho$. That total is the sum of two terms moving in opposite directions
and by very different amounts: measured in the units of
the objective~\eqref{eq:globalobj}, the travel our proposed method saves
outweighs the lateness it gives up for that saving by a factor of 7 to
29 across the sweep. The cost is therefore lower than every one of the
six strengthened rules at every penalty rate, and the gain is a large saving
in travel bought with a small concession in tardiness.

Chainability and compactness are the same property. A group whose requests chain
is one whose courier can leave each request and reach the next without a long
empty run, which is a group with short travel between its stops, and the marginal
cost of~\eqref{eq:marginal} selects for it directly. Punctuality is available to
either method once $\rho$ makes it worth buying, both driving lateness towards
zero by $\rho=10$; what differs is the distance each pays to get there. The
contribution is therefore the grouping criterion and its effect on cost, and not
an improvement in punctuality.

\subsection{The dispatch charge and the fleet}\label{sec:fleetcharge}
\titledParagraph{Behaviour without a dispatch charge.}\label{sec:phizero}
Everything above prices a dispatched courier at $\varphi = 100$. Setting
$\varphi = 0$ removes the fleet decision from the objective: any number of
couriers may be opened at no charge, and the problem reduces to minimising
travel plus lateness. Since the dispatch charge is what no published clustering
rule can respond to, this is the setting in which our advantage over them should
be most at risk, and \tabref{tab:phizero} repeats the whole comparison there.

\begin{table}[t]\centering
\sbox{\tabbox}{\small\begin{tabular}{@{}l rrrr rrrr@{}}
\toprule
& \multicolumn{4}{c}{as published} & \multicolumn{4}{c}{+ descent}\\
\cmidrule(lr){2-5}\cmidrule(lr){6-9}
method & travel & tard. & cour. & cost & travel & tard. & cour. & cost\\
\midrule
\textsc{sweep} & 2163 & 833 & 22.7 & 2996 & 1786 & 8 & 30.0 & 1794\\
\textsc{fisher} & 2487 & 578 & 34.0 & 3064 & 1605 & 5 & 32.4 & 1610\\
$k$-means & 2174 & 410 & 34.0 & 2584 & 1829 & 11 & 31.8 & 1839\\
spatio-temporal $k$-means & 2217 & 800 & 34.0 & 3017 & 1936 & 28 & 32.8 & 1964\\
capacitated $k$-means & 2284 & 449 & 34.0 & 2733 & 1868 & 10 & 32.1 & 1878\\
Ward agglomerative & 2176 & 326 & 34.0 & 2502 & 1867 & 40 & 32.1 & 1907\\
\midrule
\textbf{this paper} & \multicolumn{4}{c}{} & 1095 & 10 & 31.5 & 1106\\
ALNS & \multicolumn{4}{c}{} & 1104 & 40 & 30.8 & 1144\\
\bottomrule\end{tabular}}%
\setlength{\tabw}{0.80\textwidth}%
\ifdim\wd\tabbox>\tabw\setlength{\tabw}{\wd\tabbox}\fi
\begin{minipage}{\tabw}
\centering
\caption{$|\mathcal{R}| = 40, 100, 200$ requests, $|\mathcal{K}| = 12, 30, 60$ couriers, $Q_k = 5$, $\varphi = 0$, $\rho = 1$. The same instances with no dispatch charge, $\varphi = 0$, so cost is travel${}+{}$tardiness alone and the courier count is free. Compare with \tabref{tab:main}.}
\label{tab:phizero}
\small\begin{tabular*}{\tabw}{@{\extracolsep{\fill}}l rrrr rrrr@{}}
\toprule
& \multicolumn{4}{c}{as published} & \multicolumn{4}{c}{+ descent}\\
\cmidrule(lr){2-5}\cmidrule(lr){6-9}
method & travel & tard. & cour. & cost & travel & tard. & cour. & cost\\
\midrule
\textsc{sweep} & 2163 & 833 & 22.7 & 2996 & 1786 & 8 & 30.0 & 1794\\
\textsc{fisher} & 2487 & 578 & 34.0 & 3064 & 1605 & 5 & 32.4 & 1610\\
$k$-means & 2174 & 410 & 34.0 & 2584 & 1829 & 11 & 31.8 & 1839\\
spatio-temporal $k$-means & 2217 & 800 & 34.0 & 3017 & 1936 & 28 & 32.8 & 1964\\
capacitated $k$-means & 2284 & 449 & 34.0 & 2733 & 1868 & 10 & 32.1 & 1878\\
Ward agglomerative & 2176 & 326 & 34.0 & 2502 & 1867 & 40 & 32.1 & 1907\\
\midrule
\textbf{this paper} & \multicolumn{4}{c}{} & 1095 & 10 & 31.5 & 1106\\
ALNS & \multicolumn{4}{c}{} & 1104 & 40 & 30.8 & 1144\\
\bottomrule\end{tabular*}
\end{minipage}
\end{table}

The advantage over the clustering rules does not require the dispatch charge,
and is larger without one. Against the six rules as published our proposed method is
$-49.9\%$ (18/18 instances, $p<10^{-3}$), and
against the same rules carrying our descent $-26.5\%$
(18/18, $p<10^{-3}$), against $-12.1\%$ for the
same comparison at $\varphi = 100$. The apparent strengthening is largely
arithmetic: at $\varphi = 100$ every method pays a fixed term of $\varphi$ times
the couriers it opens, 23 on average over these instances, a large
common addend that compresses every relative difference. In absolute terms nothing changes, and the travel column shows it: our proposed
method is close to ALNS there and the rules are not.


The courier counts separate the methods by whether the fleet size is a decision
at all. The number of groups is an input to every published rule, so each of the
6 dispatches the same couriers whichever charge applies: the counts in the \emph{as published} block of \tabref{tab:main}
and of \tabref{tab:phizero} are identical. Our proposed method opens
8.3 more couriers when the charge is lifted and ALNS
7.8 more, both responding to what the objective now prices. The
strengthened rules respond as well, by 3.7 to 7.3
couriers, because the descent they carry minimises the objective, and the
objective charges $\varphi$ for every courier opened.

This is the form in which the fleet result holds. It is a claim about
responsiveness to $\varphi$ and it separates our proposed method from the
published rules, not from ALNS. It is not a claim about distance from the
capacity floor of Section~\ref{sec:problemFormulation}: \textsc{sweep} fills
each courier to $Q_k$ in angular order, so it opens
$\lceil|\mathcal{R}|/\max_k Q_k\rceil$ couriers exactly, sitting on the floor
at $\varphi=100$ and at $\varphi=0$ alike.


Against ALNS the two are not statistically distinguishable: pooled
$-0.5\%$ with 11 wins in 18, $p=0.167$. That
pooled figure hides a trend with size, $+9.6\%$ at $|\mathcal{R}|=40$, $-5.8\%$ at $|\mathcal{R}|=100$, $-5.2\%$ at $|\mathcal{R}|=200$, so ALNS is ahead on the
smallest instances and behind on the larger ones.
Under a dispatch charge the decomposition beats ALNS at matched
wall-clock and ALNS does not catch up given far longer
(Section~\ref{sec:crossover}); with couriers free it is a tie at matched
wall-clock, and a search given more time would be expected to pass us, since
what remains in the objective is exactly the route quality a search over
complete solutions is good at.

Against every published clustering rule, in its own form and strengthened with
our improvement step, our proposed method is cheaper on all 18 instances
at $\varphi = 100$ and on all 18 at $\varphi = 0$, and against ALNS it is
ahead under a dispatch charge and level without one. Both comparisons use the
three problem sizes of \tabref{tab:main}. Section~\ref{sec:scaling} takes the
one against the clustering rules to $1000$ requests and
$300$ couriers, where most of our runtime is work that is
independent per courier: our proposed method returns in
50\,s against
821\,s for the strongest clustering baseline, at a cost
34\% below
it.

\titledParagraph{The charge swept.}\label{sec:phisweep}
Section~\ref{sec:rules} prices a dispatched courier at $\varphi=100$ and
Section~\ref{sec:phizero} at $\varphi=0$. Sweeping $\varphi$ over
6 levels from $0$ to $400$ at
$|\mathcal{R}|=40$ gives the shape between them;
\figrefp{fig:sensitivity}{a} plots the result. The margin over the
strengthened rules is monotone and \emph{largest} where the charge is
smallest, $-15.1\%$ at $\varphi=0$ against
$-4.2\%$ at $\varphi=400$, which is the compression
effect of Section~\ref{sec:phizero} seen across the whole range rather than at
its ends: a term every method pays alike grows with $\varphi$ and shrinks every
ratio. Our fleet falls from 11.2
couriers at $\varphi=0$ to the capacity floor of $8$ by
$\varphi=100$ and stays there, while no rule in its published form moves
at any charge. Against ALNS the sign changes
inside the range: we are behind at $\varphi=0$ and ahead from
$\varphi=50$ upwards, which locates the crossover the earlier
sections describe.

\begin{table}[t]\centering
\sbox{\tabbox}{\small\begin{tabular}{@{}rr rrr r@{}}\toprule
& & \multicolumn{3}{c}{vs.\ rules $+$ descent} & vs.\ ALNS\\
\cmidrule(lr){3-5}\cmidrule(lr){6-6}
$|\mathcal{K}|$ & $|\mathcal{R}|$ & mean & wins & $p$ & mean\\\midrule
$6$ & $20$ & $-4.1\%$ & 6/8 & $0.055$ & $+1.7\%$\\
$9$ & $32$ & $-6.5\%$ & 7/8 & $0.039$ & $-1.1\%$\\
$12$ & $40$ & $-9.1\%$ & 8/8 & $0.008$ & $-4.3\%$\\
$18$ & $60$ & $-12.2\%$ & 8/8 & $0.008$ & $-10.7\%$\\
\bottomrule\end{tabular}}%
\setlength{\tabw}{0.80\textwidth}%
\ifdim\wd\tabbox>\tabw\setlength{\tabw}{\wd\tabbox}\fi
\begin{minipage}{\tabw}
\centering
\caption{$|\mathcal{R}| = 20, 32, 40, 60$ requests, $|\mathcal{K}| = 6, 9, 12, 18$ couriers, $Q_k = 5$, $\varphi = 100$, $\rho = 1$. The advantage widens with the fleet. The design rule of Section~\ref{sec:setup} is held fixed, so the workload grows with the fleet. Paired against the best of the six rules carrying our descent, and against ALNS at matched wall-clock, at $\varphi=100$. Negative favours the proposed method.}
\label{tab:fleet}
\small\begin{tabular*}{\tabw}{@{\extracolsep{\fill}}rr rrr r@{}}\toprule
& & \multicolumn{3}{c}{vs.\ rules $+$ descent} & vs.\ ALNS\\
\cmidrule(lr){3-5}\cmidrule(lr){6-6}
$|\mathcal{K}|$ & $|\mathcal{R}|$ & mean & wins & $p$ & mean\\\midrule
$6$ & $20$ & $-4.1\%$ & 6/8 & $0.055$ & $+1.7\%$\\
$9$ & $32$ & $-6.5\%$ & 7/8 & $0.039$ & $-1.1\%$\\
$12$ & $40$ & $-9.1\%$ & 8/8 & $0.008$ & $-4.3\%$\\
$18$ & $60$ & $-12.2\%$ & 8/8 & $0.008$ & $-10.7\%$\\
\bottomrule\end{tabular*}
\end{minipage}
\end{table}
\titledParagraph{The fleet size.}\label{sec:fleetsweep}
The design rule of Section~\ref{sec:setup} ties the request count to the
fleet, so enlarging $|\mathcal{K}|$ enlarges the instance with it and the
workload per courier is held fixed. \tabref{tab:fleet} sweeps
$|\mathcal{K}|$ from 6 to 18.

Both margins move monotonically with the fleet: against the strengthened rules
from $-4.1\%$ at $|\mathcal{K}|=6$ to
$-12.2\%$ at $|\mathcal{K}|=18$, and against
ALNS from $+1.7\%$ to
$-10.7\%$, crossing from behind to ahead. The reason is the
fleet decision itself: the number of couriers a solution may open grows with
$|\mathcal{K}|$, so there is more for the assignment program to decide and
more for a method that cannot decide it to lose.

\subsection{Scaling and parallel structure}\label{sec:scaling}
\begin{table}[t]\centering
\sbox{\tabbox}{\small\begin{tabular}{@{}r l rrr r rr@{}}\toprule
$|\mathcal{R}|$ & method & travel & tard. & cour. & cost & time & assign.\\\midrule
$100$ & ours & 1117 & 4 & 20.5 & 3171 & 7\,s & 6\,s\\
 & baseline & 1936 & 20 & 26.5 & 4606 & 1\,s & ---\\
 & ALNS & 1589 & 25 & 20.5 & 3664 & 11\,s & ---\\
\addlinespace[2pt]
$200$ & ours & 1847 & 1 & 40.5 & 5898 & 18\,s & 13\,s\\
 & baseline & 3246 & 4 & 47.5 & 8000 & 8\,s & ---\\
 & ALNS & 3235 & 59 & 40.5 & 7344 & 19\,s & ---\\
\addlinespace[2pt]
$400$ & ours & 3527 & 14 & 81.0 & 11642 & 31\,s & 20\,s\\
 & baseline & 7115 & 41 & 104.0 & 17556 & 36\,s & ---\\
 & ALNS & 6753 & 158 & 80.5 & 14961 & 90\,s & ---\\
\addlinespace[2pt]
$1000$ & ours & 7878 & 24 & 201.0 & 28002 & 50\,s & 20\,s\\
 & baseline & 16448 & 63 & 262.0 & 42711 & 821\,s & ---\\
 & ALNS & 23518 & 2602 & 200.5 & 46170 & 1483\,s & ---\\
\bottomrule\end{tabular}}%
\setlength{\tabw}{0.80\textwidth}%
\ifdim\wd\tabbox>\tabw\setlength{\tabw}{\wd\tabbox}\fi
\begin{minipage}{\tabw}
\centering
\caption{$|\mathcal{R}| = 100, 200, 400, 1000$ requests, $|\mathcal{K}| = 30, 60, 120, 300$ couriers, $Q_k = 5$, $\varphi = 100$, $\rho = 1$. Two misaligned instances at each size; those at $|\mathcal{R}| = 100$ and $200$ are also in \tabref{tab:main}, those at $400$ and $1000$ are not. Scaling of the proposed method against the strongest clustering baseline (capacitated $k$-means with our descent) on the same instances. \emph{Assign.} is the part of our time that does not divide across cores.}
\label{tab:scale}
\small\begin{tabular*}{\tabw}{@{\extracolsep{\fill}}r l rrr r rr@{}}\toprule
$|\mathcal{R}|$ & method & travel & tard. & cour. & cost & time & assign.\\\midrule
$100$ & ours & 1117 & 4 & 20.5 & 3171 & 7\,s & 6\,s\\
 & baseline & 1936 & 20 & 26.5 & 4606 & 1\,s & ---\\
 & ALNS & 1589 & 25 & 20.5 & 3664 & 11\,s & ---\\
\addlinespace[2pt]
$200$ & ours & 1847 & 1 & 40.5 & 5898 & 18\,s & 13\,s\\
 & baseline & 3246 & 4 & 47.5 & 8000 & 8\,s & ---\\
 & ALNS & 3235 & 59 & 40.5 & 7344 & 19\,s & ---\\
\addlinespace[2pt]
$400$ & ours & 3527 & 14 & 81.0 & 11642 & 31\,s & 20\,s\\
 & baseline & 7115 & 41 & 104.0 & 17556 & 36\,s & ---\\
 & ALNS & 6753 & 158 & 80.5 & 14961 & 90\,s & ---\\
\addlinespace[2pt]
$1000$ & ours & 7878 & 24 & 201.0 & 28002 & 50\,s & 20\,s\\
 & baseline & 16448 & 63 & 262.0 & 42711 & 821\,s & ---\\
 & ALNS & 23518 & 2602 & 200.5 & 46170 & 1483\,s & ---\\
\bottomrule\end{tabular*}
\end{minipage}
\end{table}

\tabref{tab:scale} runs the comparison at four sizes up
to $|\mathcal{R}| = 1000$ requests and
$|\mathcal{K}| = 300$ couriers,
five times the largest instance in
\tabref{tab:main}. At that size our proposed method returns a
solution in 50\,s against 821\,s for the strongest
clustering baseline, and it is ahead on every component rather than on the
total alone: travel 7878 against 16448,
lateness 24 against 63, and
201 couriers against 262, for a cost
34\% below the
baseline's. Of that time, 30\,s is spent in work that is independent per
courier, namely marginal costs, descent screening and exact routing, and
20\,s in the assignment program, which is not. With one core per courier
the first term divides and the second does not, so the attainable speedup is
about $2.5\times$. The
decomposition therefore delivers a genuinely parallel workload, and the
candidate list of Section~\ref{sec:assign} contributes as well, by making
every step linear in $|\mathcal{R}|$.


The ALNS rows of \tabref{tab:scale} are not a comparison at equal wall-clock,
as its time column shows. ALNS is given our runtime on the same instance as its
budget, but a budget
can only be tested between iterations, so the search runs to the end of whichever
ruin-and-recreate pass is in flight when the budget expires. At
$100$ and $200$ requests a pass is short enough that this
costs almost nothing and ALNS keeps to the budget. At $400$ it takes
three times what it was allowed, and at
$1000$, where one pass over the whole plan is itself longer
than the budget, it takes 1483\,s against the
50\,s it was allowed, a factor of
30. Even
with that much extra time its solution is far worse on the two components a
search should be good at, travel 23518 against
7878 and lateness 2602 against
24, while it opens the same fleet as we do,
200 couriers against 201. A search over complete solutions
must repair a thousand-request plan at every step, while the decomposition never
sequences more than $Q_k$ requests at once.

\subsection{Benchmark instances and the boundary of the advantage}\label{sec:external}
\titledParagraph{Adapted benchmark instances.}\label{sec:lilim}
The instances above come from our own generator, so we repeat the comparison
analysis on batches drawn from the Li \& Lim benchmark~\citep{SINTEFLiLimBenchmark}, adapted to
this setting as Section~\ref{sec:setup} describes. These have window widths,
spatial layouts and release patterns we did not choose, and they sit at the
misaligned end of the alignment axis throughout. \tabref{tab:lilim} reports
them by family.

\begin{table}[t]\centering
\sbox{\tabbox}{\small\begin{tabular}{@{}l rrrr@{}}\toprule
& ours & best of six & $+$ descent & ALNS\\\midrule
\multicolumn{5}{@{}l}{\emph{cost, by family}}\\
\quad\texttt{lc201} & 1632 & 1799 & 1663 & 1676\\
\quad\texttt{lc202} & 1581 & 1756 & 1631 & 1556\\
\quad\texttt{lr201} & 1779 & 1875 & 1814 & 1833\\
\quad\texttt{lr202} & 1977 & 2208 & 2070 & 2060\\
\quad\texttt{lrc101} & 2449 & 2827 & 2463 & 2475\\
\quad\texttt{lrc102} & 2271 & 2787 & 2258 & 2247\\
\quad\texttt{lrc201} & 2151 & 2517 & 2227 & 2161\\
\quad\texttt{lrc202} & 2125 & 2546 & 2231 & 2101\\
\addlinespace[2pt]\midrule
\multicolumn{5}{@{}l}{\emph{pooled, by component}}\\
\quad travel & 1066 & 1249 & 1084 & 1076\\
\quad tardiness & 47 & 140 & 44 & 54\\
\quad couriers & 8.8 & 9.0 & 9.2 & 8.8\\
\quad \textbf{cost} & \textbf{1996} & 2289 & 2045 & 2014\\
\bottomrule\end{tabular}}%
\setlength{\tabw}{0.80\textwidth}%
\ifdim\wd\tabbox>\tabw\setlength{\tabw}{\wd\tabbox}\fi
\begin{minipage}{\tabw}
\centering
\caption{$|\mathcal{R}| = 40$ requests, $|\mathcal{K}| = 12$ couriers, $Q_k = 5$, $\varphi = 100$, $\rho = 1$. Adapted Li \& Lim batches. \emph{Upper block}: mean cost by family. \emph{Lower block}: the same batches pooled and split into the terms of~\eqref{eq:globalobj}. \emph{Best of six} takes the cheapest published rule on each instance, \emph{+ descent} the cheapest after our improvement step.}
\label{tab:lilim}
\small\begin{tabular*}{\tabw}{@{\extracolsep{\fill}}l rrrr@{}}\toprule
& ours & best of six & $+$ descent & ALNS\\\midrule
\multicolumn{5}{@{}l}{\emph{cost, by family}}\\
\quad\texttt{lc201} & 1632 & 1799 & 1663 & 1676\\
\quad\texttt{lc202} & 1581 & 1756 & 1631 & 1556\\
\quad\texttt{lr201} & 1779 & 1875 & 1814 & 1833\\
\quad\texttt{lr202} & 1977 & 2208 & 2070 & 2060\\
\quad\texttt{lrc101} & 2449 & 2827 & 2463 & 2475\\
\quad\texttt{lrc102} & 2271 & 2787 & 2258 & 2247\\
\quad\texttt{lrc201} & 2151 & 2517 & 2227 & 2161\\
\quad\texttt{lrc202} & 2125 & 2546 & 2231 & 2101\\
\addlinespace[2pt]\midrule
\multicolumn{5}{@{}l}{\emph{pooled, by component}}\\
\quad travel & 1066 & 1249 & 1084 & 1076\\
\quad tardiness & 47 & 140 & 44 & 54\\
\quad couriers & 8.8 & 9.0 & 9.2 & 8.8\\
\quad \textbf{cost} & \textbf{1996} & 2289 & 2045 & 2014\\
\bottomrule\end{tabular*}
\end{minipage}
\end{table}

Pooled over the 24 benchmark batches our proposed method is
$-12.2\%$ against the best published rule (24 wins,
$p<10^{-3}$), a figure that coincides to three significant figures with an
unrelated one, the $-12.2\%$ against \emph{strengthened}
rules at $|\mathcal{K}|=18$ in \tabref{tab:fleet}, which is
measured on generated instances; the two are different comparisons on different
data and agree by chance. Against the best of those rules carrying our descent
the margin here is $-2.4\%$
(18 wins, $p=0.001$). Against ALNS at matched
wall-clock it is $-0.9\%$ (15/24,
$p=0.101$), a tie at this number of instances.
The lower block of \tabref{tab:lilim} gives the components, which locate
that margin: travel and lateness are nearly level with the best strengthened rule
and with ALNS, and most of the difference comes from the fleet size.

Both figures are measured at $|\mathcal{R}| = 40$:
$-8.9\%$ on our generated instances against
$-2.4\%$ on the benchmark batches. Instance size is therefore not
what makes the margin smaller here.
Section~\ref{sec:instancespace} identifies it.

\titledParagraph{Instance structure and the size of the advantage.}\label{sec:instancespace}
The margin on the benchmark batches is smaller than on the instances generated
by the design rule of Section~\ref{sec:setup}. Whether that rule yields
conditions favourable to the proposed method is a question about instance
structure, and we answer it by moving the generator along the axes on which the
two datasets differ and watching the margin. Three axes were tried.
\tabref{tab:instancespace} carries the two swept here; the
alignment axis was swept already in Section~\ref{sec:alignaxis}.
Two do not explain the gap and one
does. The two failures are reported because they say what the effect is
\emph{not}.

\begin{table}[t]\centering
\sbox{\tabbox}{\small\begin{tabular}{@{}l r@{\;\;} r@{\;\;} r@{}}\toprule
setting & property & $+$ desc. & vs.\ ALNS\\\midrule
\multicolumn{4}{@{}l}{\emph{window width: courier-time ratio}}\\
\quad width $45$ & $1.77$ & $-9.1\%$ & $-4.1\%$\\
\quad width $90$ & $2.05$ & $-13.4\%$ & $-5.9\%$\\
\quad width $180$ & $2.61$ & $-14.9\%$ & $-4.2\%$\\
\quad width $360$ & $3.75$ & $-14.3\%$ & $-3.6\%$\\
\quad width $720$ & $6.01$ & $-14.9\%$ & $-4.1\%$\\
\addlinespace[2pt]
\multicolumn{4}{@{}l}{\emph{delivery placement: locality}}\\
\quad local & $0.85$ & $-9.1\%$ & $-3.6\%$\\
\quad cross & $0.03$ & $-2.0\%$ & $-0.2\%$\\
\midrule
benchmark (Li \& Lim) & $0.29$ & $-2.4\%$ & $-0.9\%$\\
\bottomrule\end{tabular}}%
\setlength{\tabw}{0.80\textwidth}%
\ifdim\wd\tabbox>\tabw\setlength{\tabw}{\wd\tabbox}\fi
\begin{minipage}{\tabw}
\centering
\caption{$|\mathcal{R}| = 40$ requests, $|\mathcal{K}| = 12$ couriers, $Q_k = 5$, $\varphi = 100$, $\rho = 1$. Moving our generator towards the benchmark, one axis at a time. Each block reports the measured instance property and the margin against the best of the six rules carrying our descent, and against ALNS. The benchmark row is the adapted Li \& Lim batches at the same size.}
\label{tab:instancespace}
\small\begin{tabular*}{\tabw}{@{\extracolsep{\fill}}l r@{\;\;} r@{\;\;} r@{}}\toprule
setting & property & $+$ desc. & vs.\ ALNS\\\midrule
\multicolumn{4}{@{}l}{\emph{window width: courier-time ratio}}\\
\quad width $45$ & $1.77$ & $-9.1\%$ & $-4.1\%$\\
\quad width $90$ & $2.05$ & $-13.4\%$ & $-5.9\%$\\
\quad width $180$ & $2.61$ & $-14.9\%$ & $-4.2\%$\\
\quad width $360$ & $3.75$ & $-14.3\%$ & $-3.6\%$\\
\quad width $720$ & $6.01$ & $-14.9\%$ & $-4.1\%$\\
\addlinespace[2pt]
\multicolumn{4}{@{}l}{\emph{delivery placement: locality}}\\
\quad local & $0.85$ & $-9.1\%$ & $-3.6\%$\\
\quad cross & $0.03$ & $-2.0\%$ & $-0.2\%$\\
\midrule
benchmark (Li \& Lim) & $0.29$ & $-2.4\%$ & $-0.9\%$\\
\bottomrule\end{tabular*}
\end{minipage}
\end{table}

\titledParagraph{Not the alignment.} Section~\ref{sec:alignaxis} places the
benchmark at a realised misalignment of $0.61$, beside our hardest setting,
where our margin is at its largest. If anything the benchmark sits where the
method should do best.

\titledParagraph{Not the time pressure.} Widening the time windows moves the
courier-time ratio from $1.77$ to
$6.01$, across and beyond the benchmark's $2.61$. The
margin does not shrink; it \emph{grows}, from $-9.1\%$ to
$-14.9\%$. At the width whose courier-time ratio matches the
benchmark exactly the margin is $-14.9\%$,
six times the benchmark's. How much time the fleet has is not the explanation,
and we had expected it to be.

\titledParagraph{The spatial structure.} What does explain it is where a
request's delivery lies. Our default generator puts the delivery near its own
pickup, so the locality defined in Section~\ref{sec:setup}, the fraction of
requests whose pickup and delivery lie within $10$ of each other, is
$0.85$. Sending the delivery to a different hub gives
$0.03$, and the benchmark sits between the two at $0.29$. The margin
against the strengthened rules collapses from $-9.1\%$ to
$-2.0\%$ (7/8, $p=0.109$), and against
ALNS from $-3.6\%$ to $-0.2\%$, within noise of
the $-2.4\%$ and $-0.9\%$ measured on the benchmark
itself.

The mechanism is the one Section~\ref{sec:cost} argues for, seen at its
boundary. When a
delivery sits near its pickup, a request is effectively a point, the freedom
lies in which points travel together, and what a group costs is exactly the
chaining the marginal cost measures and a centroid cannot. When the delivery is
across the city, every request is a long mandatory traversal, the route is
dominated by edges no grouping can avoid, and the choice of partition matters
much less: any sensible grouping followed by a good descent reaches a similar
place. The benchmark lies towards that end, which is why our advantage there is
real but small.

Courier dispatch sits at the other end of the same axis. In meal delivery, the
setting whose operational conditions this paper
takes as its model~\citep{reyes2018meal, yildiz2019provably,
ulmer2021restaurant}, and in comparable work such as pharmacy delivery, a
customer orders from a vendor already close by and the platform serves the order
within a bounded radius, so a request's pickup and its delivery are near each
other. That is why the generator of Section~\ref{sec:setup} places each
delivery near its own pickup, and it is the regime in which the grouping
criterion rather than the unavoidable edges decides what a solution costs. The
boundary is stated here because the margin is not uniform across instance
structures, and because the structures in which it is largest are the ones the
application presents.

\subsection{Comparison with the metaheuristic reference}

\tabref{tab:main} also lists an adaptive large neighbourhood search under the
budget of Section~\ref{sec:setup}, our own wall-clock on each instance with a
ten-second floor. Under that budget
the decomposition is $-14.2\%$ against it, winning 18 of
18 instances ($p<10^{-3}$), and \tabref{tab:bysize} splits that by
size. The difference is not the fleet, since the two dispatch within
3.7\% of each other at every size. It is travel, where our margin
widens as the instances grow: $4\%$ at $|\mathcal{R}|=40$, $37\%$ at $|\mathcal{R}|=100$, $47\%$ at $|\mathcal{R}|=200$.
The explanation is the budget rather than the algorithm. A ruin-and-recreate
step reinserts requests into a plan spanning the whole instance, so the work per
improvement grows with $|\mathcal{R}|$, while the decomposition never sequences
more than $Q_k$ requests at a time and solves each of those exactly. Given our
wall-clock, ALNS has fewer useful moves the larger the instance is.
Section~\ref{sec:crossover} removes the budget restriction and measures how
much more time it needs.

\titledParagraph{Time for the reference to match.}\label{sec:crossover}
\begin{table}[t]\centering
\sbox{\tabbox}{\small\begin{tabular}{@{}r@{\;\;} l@{\;\;} r@{\;\;}r@{\;\;}r@{\;\;} r@{\;\;} r@{}}\toprule
$|\mathcal{R}|$ & & travel & tard. & cour. & cost & time\\\midrule
$100$ & ours & 1117 & 4 & 20.5 & 3171 & 7.6\,s\\
 & ALNS & 1252 & 6 & 20.5 & 3307 ($+4\%$) & 300\,s\\
\addlinespace[2pt]
$200$ & ours & 1847 & 1 & 40.5 & 5898 & 18.1\,s\\
 & ALNS & 2454 & 25 & 41.5 & 6629 ($+12\%$) & 300\,s\\
\addlinespace[2pt]
$400$ & ours & 3527 & 14 & 81.0 & 11642 & 37.7\,s\\
 & ALNS & 6129 & 117 & 81.0 & 14345 ($+23\%$) & 300\,s\\
\bottomrule\end{tabular}}%
\setlength{\tabw}{0.80\textwidth}%
\ifdim\wd\tabbox>\tabw\setlength{\tabw}{\wd\tabbox}\fi
\begin{minipage}{\tabw}
\centering
\caption{$|\mathcal{R}| = 100, 200, 400$ requests, $|\mathcal{K}| = 30, 60, 120$ couriers, $Q_k = 5$, $\varphi = 100$, $\rho = 1$. ALNS is stopped as soon as its incumbent reaches the cost of the proposed method, or after 300\,s. It never reached it. Components are shown for both, so the residual gap can be attributed. Two misaligned instances per size, each searched by ALNS with two seeds; the ALNS row is the worse of the two seeds on each instance, the only one whose components were recorded. The better seed is $+4\%$, $+9\%$ and $+21\%$ above the proposed method at the three sizes.}
\label{tab:crossover}
\small\begin{tabular*}{\tabw}{@{\extracolsep{\fill}}r@{\;\;} l@{\;\;} r@{\;\;}r@{\;\;}r@{\;\;} r@{\;\;} r@{}}\toprule
$|\mathcal{R}|$ & & travel & tard. & cour. & cost & time\\\midrule
$100$ & ours & 1117 & 4 & 20.5 & 3171 & 7.6\,s\\
 & ALNS & 1252 & 6 & 20.5 & 3307 ($+4\%$) & 300\,s\\
\addlinespace[2pt]
$200$ & ours & 1847 & 1 & 40.5 & 5898 & 18.1\,s\\
 & ALNS & 2454 & 25 & 41.5 & 6629 ($+12\%$) & 300\,s\\
\addlinespace[2pt]
$400$ & ours & 3527 & 14 & 81.0 & 11642 & 37.7\,s\\
 & ALNS & 6129 & 117 & 81.0 & 14345 ($+23\%$) & 300\,s\\
\bottomrule\end{tabular*}
\end{minipage}
\end{table}

Given the same wall-clock the proposed method is ahead, and the natural objection
is that ALNS would win with more time. \tabref{tab:crossover} tests
that directly, by stopping the search the moment its incumbent reaches our
cost. Over every instance and both seeds it never did, inside a budget of
300\,s against our 7--38\,s, and the shortfall
widens with size, reaching $+23\%$ at
$|\mathcal{R}| = 400$, and $+21\%$ for the better of the two seeds.


The components say what the residual is, and the obvious guess about it is
wrong. One might expect a ruin-and-recreate search to be held
back by the fleet. Closing a courier means relocating every one of its requests,
and every intermediate step raises the objective, because a route with one
request left still pays $\varphi$ in full. No standard removal rule performs that
transformation in a single move. That is not what happens. At $|\mathcal{R}| = 400$ ALNS dispatches
81.0 couriers against our
81.0, so it sizes the fleet as we do. The gap is
travel, 2601 units of it, about
96\% of the shortfall.

The reason is the size of the neighbourhood rather than its shape, and it is the
one given above: each ruin-and-recreate step works on a plan spanning the whole
instance, so three hundred seconds buys proportionally fewer useful moves the
larger that instance is, which is why the shortfall widens rather than closes. We would expect a longer budget to close it eventually; what
we can report is that it does not close within the $300$\,s cap, which is eight
to forty-four times our own runtime depending on the instance.


Two caveats keep this in proportion. First, the effect depends on
$\varphi > 0$: Section~\ref{sec:phizero} repeats the whole comparison with
dispatch free, and there the two are indistinguishable at matched wall-clock, so
ALNS would be expected to pass us given more time. Second, the
comparison uses the standard operator set of~\citet{ropke2006adaptive}; an
implementation tuned for these sizes, or one carrying a route-elimination
operator, would be expected to do better than this.

\subsection{The parallel assignment}\label{sec:lagrangianresults}

Section~\ref{sec:lagrangian} replaces the one coupled program by one
subproblem per courier. \tabref{tab:lagrangian} reports what that costs, in
components rather than a single percentage, on the same 18 instances as
\tabref{tab:main}.

\begin{table}[t]\centering
\sbox{\tabbox}{\small\begin{tabular}{@{}r@{\;\;} l@{\;} r@{\;\;}r@{\;\;}r@{\;\;} r@{\;\;} r@{\;\;} r@{}}\toprule
$|\mathcal{R}|$ & assignment & travel & tard. & cour. & cost & assign. & $\Delta$\\\midrule
$40$ & exact program & 624 & 18 & 8.0 & 1441 & 0.2\,s & ---\\
 & multipliers only & 547 & 6 & 10.0 & 1553 & 0.1\,s & $+7.8\%$\\
 & $+$ fleet refit & 593 & 13 & 8.7 & 1472 & 0.1\,s & $+2.2\%$\\
 & $+$ shortlist refit & 614 & 10 & 8.3 & 1457 & 0.3\,s & $+1.1\%$\\
\addlinespace[2pt]
$100$ & exact program & 1121 & 67 & 20.7 & 3255 & 3.3\,s & ---\\
 & multipliers only & 1078 & 11 & 24.5 & 3539 & 0.3\,s & $+8.7\%$\\
 & $+$ fleet refit & 1057 & 17 & 23.7 & 3441 & 0.3\,s & $+5.7\%$\\
 & $+$ shortlist refit & 1096 & 81 & 21.2 & 3294 & 1.5\,s & $+1.2\%$\\
\addlinespace[2pt]
$200$ & exact program & 1874 & 231 & 40.8 & 6189 & 10.5\,s & ---\\
 & multipliers only & 1830 & 34 & 46.5 & 6515 & 0.5\,s & $+5.3\%$\\
 & $+$ fleet refit & 1850 & 72 & 46.0 & 6522 & 0.6\,s & $+5.5\%$\\
 & $+$ shortlist refit & 1870 & 204 & 41.2 & 6191 & 6.6\,s & $+0.1\%$\\
\addlinespace[2pt]
\midrule
\emph{all} & exact program & 1206 & 105 & 23.2 & 3628 & 4.7\,s & ---\\
 & multipliers only & 1152 & 17 & 27.0 & 3869 & 0.3\,s & $+7.3\%$\\
 & $+$ fleet refit & 1166 & 34 & 26.1 & 3812 & 0.3\,s & $+4.5\%$\\
 & $+$ shortlist refit & 1193 & 98 & 23.6 & 3647 & 2.8\,s & $+0.8\%$\\
\bottomrule\end{tabular}}%
\setlength{\tabw}{0.80\textwidth}%
\ifdim\wd\tabbox>\tabw\setlength{\tabw}{\wd\tabbox}\fi
\begin{minipage}{\tabw}
\centering
\caption{$|\mathcal{R}| = 40, 100, 200$ requests, $|\mathcal{K}| = 12, 30, 60$ couriers, $Q_k = 5$, $\varphi = 100$, $\rho = 1$. The cost of distributing the assignment, on the instances of \tabref{tab:main} at $\varphi = 100$. \emph{Assign.} is the time in the assignment step, the only step that couples couriers; $\Delta$ is the paired cost difference against the exact program. Each block is one problem size; \emph{all} pools the instances of every size, which is what the text quotes.}
\label{tab:lagrangian}
\small\begin{tabular*}{\tabw}{@{\extracolsep{\fill}}r@{\;\;} l@{\;} r@{\;\;}r@{\;\;}r@{\;\;} r@{\;\;} r@{\;\;} r@{}}\toprule
$|\mathcal{R}|$ & assignment & travel & tard. & cour. & cost & assign. & $\Delta$\\\midrule
$40$ & exact program & 624 & 18 & 8.0 & 1441 & 0.2\,s & ---\\
 & multipliers only & 547 & 6 & 10.0 & 1553 & 0.1\,s & $+7.8\%$\\
 & $+$ fleet refit & 593 & 13 & 8.7 & 1472 & 0.1\,s & $+2.2\%$\\
 & $+$ shortlist refit & 614 & 10 & 8.3 & 1457 & 0.3\,s & $+1.1\%$\\
\addlinespace[2pt]
$100$ & exact program & 1121 & 67 & 20.7 & 3255 & 3.3\,s & ---\\
 & multipliers only & 1078 & 11 & 24.5 & 3539 & 0.3\,s & $+8.7\%$\\
 & $+$ fleet refit & 1057 & 17 & 23.7 & 3441 & 0.3\,s & $+5.7\%$\\
 & $+$ shortlist refit & 1096 & 81 & 21.2 & 3294 & 1.5\,s & $+1.2\%$\\
\addlinespace[2pt]
$200$ & exact program & 1874 & 231 & 40.8 & 6189 & 10.5\,s & ---\\
 & multipliers only & 1830 & 34 & 46.5 & 6515 & 0.5\,s & $+5.3\%$\\
 & $+$ fleet refit & 1850 & 72 & 46.0 & 6522 & 0.6\,s & $+5.5\%$\\
 & $+$ shortlist refit & 1870 & 204 & 41.2 & 6191 & 6.6\,s & $+0.1\%$\\
\addlinespace[2pt]
\midrule
\emph{all} & exact program & 1206 & 105 & 23.2 & 3628 & 4.7\,s & ---\\
 & multipliers only & 1152 & 17 & 27.0 & 3869 & 0.3\,s & $+7.3\%$\\
 & $+$ fleet refit & 1166 & 34 & 26.1 & 3812 & 0.3\,s & $+4.5\%$\\
 & $+$ shortlist refit & 1193 & 98 & 23.6 & 3647 & 2.8\,s & $+0.8\%$\\
\bottomrule\end{tabular*}
\end{minipage}
\end{table}

The multipliers $\lambda$ of~\eqref{eq:lagr} coordinate one of the two
decisions well and the other badly: which courier takes a request, and how many
couriers open. The components separate
them. The Lagrangian variant travels \emph{less} than the exact program
(1152 against 1206) and is \emph{more}
punctual (17 against 105); it loses
only on the fleet, opening 27.0 couriers against
23.2, and that one term turns a better route plan into a
$+7.3\%$ worse solution (0/18 wins,
$p<10^{-3}$). This is the expected behaviour of a Lagrangian relaxation of a
fixed-charge problem: a multiplier on the covering constraint tells a courier
which requests are cheap for it, and nothing tells the couriers collectively how
many of them should open.


Two repairs take the multipliers as a starting point rather than as the answer,
and both appear in \tabref{tab:lagrangian}. The \emph{fleet refit} keeps the
set of couriers the multipliers opened, together with the two sizes either side
of it, and assigns requests to each of those fleets in turn; with the open set
fixed the dispatch charge is a constant and what remains is a transportation
problem, so each trial is a linear program. The \emph{shortlist refit} leaves
the fleet decision open and uses the reduced costs only to narrow the problem,
keeping for each request the three couriers of lowest reduced cost and solving
the exact assignment program on what is left.

The second recovers most of the loss, $+0.8\%$ pooled, and more of
it the larger the instance:
$+1.1\%$ at $|\mathcal{R}|=40$ against
$+0.1\%$ at $|\mathcal{R}|=200$. That confirms which
term carries the gap rather than removing it. That repair is not free: the
program it solves is smaller than the original but still an integer program, and
still coupled, which costs
2.77\,s against the exact program's
4.66\,s pooled, but only because it is cheaper at the two
larger sizes; at $|\mathcal{R}|=40$, where the exact program
closes in 0.17\,s, the repair is the dearer of the two at
0.26\,s. What the distributed form buys is the coupled step:
16 times less serial time
pooled, and a saving that grows with the instance, from
1.5 times at $|\mathcal{R}|=40$ (0.17\,s against 0.11\,s) to
21 times at $|\mathcal{R}|=200$, in a step that
needs no information about other couriers. We therefore use the
exact assignment in the production method, and report the Lagrangian variant as
the option for a setting where the assignment must itself be distributed.

